% Capitulo: Governança de dados na era da IA
\chapter{Governança de dados na era da IA}
\label{cap:governanca}

Governança de dados sempre teve fama de burocracia: comitês, políticas longas,
dicionários de dados em planilha e um responsável cobrando todo mundo por cadastros mal
preenchidos. Não por acaso, era a primeira coisa cortada quando um projeto apertava. Este
capítulo argumenta que a IA muda essa equação. Boa parte do trabalho braçal da governança
(documentar, classificar, conciliar, avisar) pode ser proposta pela máquina, e o papel de
quem cuida dos dados passa de executor a tutor e revisor. E construir uma ferramenta interna
que faça parte desse trabalho ficou ao alcance de qualquer time.

\section{O que é governança de dados}
\label{sec:gov-o-que-e}

Governança de dados é o conjunto de papéis, regras e processos que garantem que os dados
são confiáveis, bem descritos, protegidos e usados de forma adequada \parencite{dama2017}.
Ela abrange várias frentes que já apareceram neste livro: o catálogo e os metadados, que
dizem o que cada dado é; a qualidade, que diz se ele está certo para o uso
(\veja{sec:fund-qualidade}); o master data (capítulo~\ref{cap:master-data}); o acesso e
a privacidade; a linhagem, que mostra de onde cada dado veio; e o ciclo de vida, que define
quanto tempo cada dado fica guardado.

Os papéis clássicos são três, embora a nomenclatura varie entre autores e empresas. O
\emph{dono do dado} é alguém do negócio que responde por um domínio e decide as suas regras.
O \emph{data steward} cuida do dia a dia: a qualidade, a documentação, a resolução de
problemas. O \emph{custodiante} cuida da parte técnica: onde o dado fica, como é protegido,
como é copiado. \pendente{conferir a nomenclatura dos papéis no DMBOK}
\textcite{eryurek2021} mostram como esses papéis se reorganizam quando os dados vão para a
nuvem, com muito mais automação e muito mais gente com acesso.

\section{Como a governança era vista}
\label{sec:gov-antes}

Na forma tradicional, governança era um programa corporativo com comitês, políticas
extensas e um conjunto de documentos que ninguém lia. O dicionário de dados ficava numa
planilha, desatualizado no dia seguinte à publicação. E o data steward fazia papel de
fiscal: mantinha tudo à mão, cobrava dos outros e estava sempre atrás dos dados, que mudavam
mais rápido do que ele conseguia documentar.

O resultado era previsível. A governança virava sinônimo de atraso, as áreas aprendiam a
contorná-la, e, no primeiro aperto de prazo, era cortada. Não por má vontade, mas porque o
custo era visível e imediato, e o benefício, difuso e futuro, como acontecia com o master
data (\veja{sec:master-data-antes}).

Uma das críticas mais úteis a esse modelo é a da governança não invasiva. Em vez de criar
uma estrutura nova, com cargos e comitês novos, a proposta é reconhecer e formalizar as
responsabilidades que as pessoas já exercem sobre os dados no trabalho do dia a dia
\parencite{seiner2014}. Quem cadastra clientes já é, na prática, responsável pela qualidade
desse cadastro. A governança deixa de ser algo imposto de fora e passa a ser a organização
do que já acontece.

\section{O data steward como tutor e revisor da IA}
\label{sec:steward-tutor}

A IA não muda o que a governança precisa fazer. Muda quem faz o trabalho braçal. Uma lista
do que ela já pode propor, com boa qualidade, num ambiente de dados:

\begin{itemize}
  \item descrições de tabelas e colunas a partir dos nomes, dos dados e do código que as gera;
  \item classificação de dados sensíveis, como dado pessoal ou financeiro, para orientar
    permissões;
  \item pares suspeitos de duplicidade em cadastros mestres;
  \item inconsistências entre fontes, como o mesmo produto com unidades diferentes;
  \item regras de qualidade a partir do comportamento dos dados, como faixas de valores
    esperadas.
\end{itemize}

Em todos esses casos, a IA propõe e uma pessoa decide. O data steward deixa de escrever
cada descrição e passa a revisar descrições propostas; deixa de procurar duplicidades e
passa a aprovar ou rejeitar pares sugeridos. E, a cada decisão, ensina: a correção vira
exemplo ou regra, e a próxima proposta vem melhor (Figura~\ref{fig:governanca-ciclo}). O
steward passa a ser tutor e revisor da IA.

\begin{figure}[htbp]
  \centering
  % Ciclo de governanca com o data steward como tutor e revisor da IA (cap. 14).
\begin{tikzpicture}[diagrama, x=1mm, y=1mm,
    etapa/.style={caixa, text width=34mm, minimum height=13mm, font=\scriptsize}]
  \node[destaquecaixa, etapa] (prop) at (55,28)
    {\textbf{A IA propõe}\\descrições, classificação de dado sensível, pares suspeitos};
  \node[etapa] (rev) at (96,10)
    {\textbf{O steward revisa}\\aprova, corrige ou rejeita};
  \node[etapa] (reg) at (55,-8)
    {\textbf{A decisão é registrada}\\auditoria, correspondência, catálogo};
  \node[etapa] (apr) at (14,10)
    {\textbf{Regras e exemplos melhoram}\\cada correção ensina a próxima proposta};
  \draw[seta] (prop.east) to[bend left=20] (rev.north);
  \draw[seta] (rev.south) to[bend left=20] (reg.east);
  \draw[seta] (reg.west) to[bend left=20] (apr.south);
  \draw[seta] (apr.north) to[bend left=20] (prop.west);
  \node[rotulo] at (55,10) {o steward como\\tutor e revisor};
\end{tikzpicture}

  \caption{O ciclo de governança com a IA propondo e o data steward revisando. Cada
    decisão registrada melhora a próxima proposta.}
  \label{fig:governanca-ciclo}
  \fonte{elaborado pelo autor.}
\end{figure}

É o mesmo movimento que aparece em outros pontos do livro. Testes escritos antes do modelo
eram recomendados e raramente feitos, e a IA os tornou baratos o bastante para virar regra
(\veja{sec:tdd}). A documentação era boa prática que sempre ficava para depois, e virou o
insumo principal dos assistentes que respondem perguntas sobre os dados
(\veja{sec:metadados-contexto}). Com a governança acontece o mesmo: o que antes não cabia no
tempo de ninguém agora cabe, desde que alguém revise.

A mudança tem uma consequência para a carreira de quem cuida dos dados. O valor do data
steward deixa de estar no volume de trabalho que ele executa e passa a estar no julgamento
que ele aplica. Conhecer o negócio, saber quando duas coisas parecidas são diferentes,
perceber que uma descrição tecnicamente correta vai confundir quem lê: é isso que a IA não
faz, e é isso que o steward ensina a ela.

\rascunho{
  Relato: onde eu já vi ou apliquei esse modelo de steward como revisor da IA (por exemplo,
  descrições de tabelas geradas e revisadas, ou classificação de dados sensíveis), e o que
  foi preciso corrigir. Sem nomes.
}

\section{Uma aplicação interna ao alcance de qualquer time}
\label{sec:gov-aplicacao-interna}

Um exemplo concreto mostra como o cenário mudou. Imagine um programa simples, em Python, de
uso interno, que roda depois de cada carga de dados. Ele procura, nas dimensões do armazém,
valores diferentes que provavelmente representam a mesma coisa: o mesmo cliente escrito de
três jeitos, o mesmo produto com duas descrições. Quando encontra um caso, registra numa
fila de revisão e manda uma mensagem para o dono do domínio, com o link para decidir.

Há poucos anos, isso seria um projeto, com especificação, orçamento e meses de
desenvolvimento. Hoje, alguém que programa em Python monta uma primeira versão em poucos
dias, com a IA ajudando a escrever o código. O que continua difícil não é o código: é
combinar com o negócio quem decide, o que conta como \enquote{o mesmo cliente} e o que
acontece depois da decisão. O capítulo~\ref{cap:aplicacao-propria} mostra o que considerar
para levar uma aplicação dessas a sério, com um trecho de código do motor de conciliação
(\veja{sec:motor-conciliacao}).

\section{Ferramentas de mercado}
\label{sec:gov-ferramentas}

Construir não é a única opção. O mercado de governança, catálogo e master data é grande, e
a Tabela~\ref{tab:ferramentas-governanca} organiza as alternativas por categoria. É um
retrato de 2026: nomes, donos e preços mudam rápido, e por isso o que importa são as
categorias, não os nomes.

\begin{table}[htbp]
  \centering
  \caption{Ferramentas de governança, catálogo e master data por categoria.}
  \label{tab:ferramentas-governanca}
  \small
  \begin{tabularx}{\textwidth}{@{}>{\raggedright\arraybackslash}p{0.25\textwidth}
      >{\raggedright\arraybackslash}X>{\raggedright\arraybackslash}X@{}}
    \toprule
    Categoria & Exemplos & Quando faz sentido \\
    \midrule
    Catálogo e governança independentes & Alation, Collibra, Atlan, Informatica (Salesforce) & Muitas plataformas de dados, setor regulado, programa corporativo \\
    \addlinespace
    Nativas da plataforma & Databricks Unity Catalog, Snowflake Horizon, Microsoft Purview & Dados concentrados numa plataforma só \\
    \addlinespace
    Código aberto & DataHub, OpenMetadata & Time com engenharia para hospedar e manter \\
    \addlinespace
    Master data & Informatica MDM (Salesforce), Reltio (SAP), Profisee, Semarchy, Stibo Systems & Vários domínios, muitas fontes, regras complexas de conciliação \\
    \addlinespace
    Aplicação própria & Python, um banco de dados e um canal de mensagens & Começar pequeno, um domínio por vez, regras muito específicas do negócio \\
    \bottomrule
  \end{tabularx}
  \fonte{elaborado pelo autor com base nos sites dos fornecedores, em novembro de 2026.
    \pendente{conferir a tabela antes da impressão}}
\end{table}

O mercado também se movimenta. Nos últimos anos, dois dos maiores fornecedores
independentes de master data foram comprados por gigantes de software corporativo: a
Informatica pela Salesforce e a Reltio pela SAP. \pendente{conferir datas e citar fontes
oficiais antes da impressão} Para quem escolhe uma ferramenta, isso reforça a importância
de avaliar não só os recursos, mas também a facilidade de tirar os dados de lá se um dia
for preciso.

Uma orientação geral ajuda a navegar a tabela. As ferramentas nativas da plataforma
resolvem boa parte dos casos de quem concentra os dados num único armazém, com pouco custo
adicional. Os catálogos independentes fazem mais sentido quando há muitas plataformas, ou em
setores regulados que exigem trilhas de auditoria formais. As opções de código aberto são
boas para quem tem engenharia para hospedar e manter. E nada impede começar com uma
aplicação própria, pequena, num domínio, e migrar quando a escala pedir.

\section{O que não se delega à IA}
\label{sec:gov-nao-delegar}

Com tanto trabalho podendo ser proposto pela máquina, vale deixar claro o que continua com
as pessoas, por definição.

\begin{itemize}
  \item \textbf{A responsabilidade por cada domínio.} Alguém precisa responder pelo cadastro
    de clientes, e esse alguém não pode ser um modelo.
  \item \textbf{As regras de negócio.} O que é um cliente ativo, quando dois cadastros
    representam a mesma empresa, qual atributo é obrigatório. A IA pode sugerir; o negócio
    decide.
  \item \textbf{O acesso e as obrigações legais.} Quem pode ver o quê, e como a empresa
    cumpre a legislação de proteção de dados pessoais \parencite{lgpd2018}.
  \item \textbf{A decisão final em caso de dúvida.} A IA propõe, uma pessoa decide e assina.
\end{itemize}

A governança na era da IA, portanto, não é menor. É diferente: menos tempo executando, mais
tempo decidindo e ensinando. E, pela primeira vez, ela tem chance de andar na mesma
velocidade que os dados.

\begin{leituras}
  \leitura{eryurek2021}{Uma visão atual e prática de governança, com atenção a dados na
    nuvem, qualidade, privacidade e ao papel de cada pessoa.}
  \leitura{seiner2014}{A ideia de governança que formaliza o que as pessoas já fazem, em
    vez de criar uma burocracia nova. Combina bem com o steward como revisor.}
  \leitura{dama2017}{O guia de referência para o vocabulário e para as áreas da gestão de
    dados.}
\end{leituras}
